% Local: run `make poster` or pdfLaTeX twice.
% Overleaf: select poster.tex as the main document and use pdfLaTeX.
\documentclass[11pt]{beamer}
\usepackage[orientation=portrait,size=a0,scale=1.32]{beamerposter}
\usetheme{TemplatePoster}

\title{Template: a research poster example}
\author{Your name \and Your collaborator}
\institute{Your research team --- Your university or organisation}
\date{Conference name and date}

\titlegraphic{%
  \logobar[0.04\paperheight]
    {assets/logos/concordia_logo.pdf}
    {assets/logos/gina_cody-wordmark.pdf}
    {assets/logos/conference_logo.pdf}%
}

\begin{document}
\begin{frame}[fragile,t]
  \begin{colorblock}{1. Context and research gap}
    Reliable comparisons depend on clear questions, consistent
    measurements, and a shared view of the study. This example follows
    a simple workflow from data collection to evaluation.
    \begin{figure}
      \figureplaceholder[0.26\linewidth]{Study overview}
      \caption{An overview of the research context and workflow.}
    \end{figure}
  \end{colorblock}

  \begin{colorblock}{2. Research questions}
    How can a simple approach improve the quality and efficiency of a
    research workflow? The study has three goals:
    \begin{itemize}
      \item Establish a reproducible baseline.
      \item Compare the proposed approach under the same conditions.
      \item Identify limitations and opportunities for further study.
    \end{itemize}
  \end{colorblock}

  \begin{colorblock}{3. Method}
    This Java method computes the mean of a set of values.
\begin{lstlisting}[language=Java,xleftmargin=2em]
public static double mean(double[] values) {
    if (values.length == 0) {
        throw new IllegalArgumentException("No values");
    }
    double sum = 0.0;
    for (double value : values) {
        sum += value;
    }
    return sum / values.length;
}
\end{lstlisting}
  \end{colorblock}

  \begin{colorblock*}{4. Visual overview}
    \begin{figure}
      \figureplaceholder[0.075\paperheight]{Main research figure}
      \caption{An overview of the research workflow.}
    \end{figure}
  \end{colorblock*}

  \begin{colorblock}{5. Comparison}
    Two views show how the baseline and proposed approach differ.
    Both use the same dataset and evaluation conditions.
    \begin{figure}
      \begin{minipage}[t]{0.48\linewidth}
        \figureplaceholder[0.65\linewidth]{Baseline}
        \caption{The reference approach.}
      \end{minipage}\hfill
      \begin{minipage}[t]{0.48\linewidth}
        \figureplaceholder[0.65\linewidth]{Proposed approach}
        \caption{The proposed approach.}
      \end{minipage}
    \end{figure}
  \end{colorblock}

  \begin{colorblock}{6. Results}
    The proposed approach achieves a higher score with a shorter run
    time in this illustrative comparison.
    \begin{table}
      \centering
      \begin{tabular}{@{}lrr@{}}
        \toprule
        \textbf{Method} & \textbf{Score} & \textbf{Time (s)} \\
        \midrule
        Baseline & 0.72 & 12.4 \\
        Variant A & 0.78 & 11.6 \\
        Variant B & 0.81 & 10.8 \\
        \textbf{Proposed} & \textbf{0.86} & \textbf{9.7} \\
        \bottomrule
      \end{tabular}
      \caption{Illustrative data only; higher scores and shorter times are better.}
    \end{table}
  \end{colorblock}

  \begin{colorblock}{7. Key findings}
    \textbf{Main finding.} The illustrative comparison shows
    \textcolor{turquoise}{\textbf{higher scores}} alongside shorter run
    times. Quality and efficiency can be considered together.
    \par\medskip
    \textbf{Why it matters.} Meaningful comparisons use
    \emph{the same data, budget, and evaluation rules}. A reproducible
    baseline makes differences easier to interpret.
    \par\medskip
    \textbf{Next step.} Evaluate the approach on
    \alert{larger, independent datasets} and report uncertainty
    before drawing broader conclusions.
  \end{colorblock}

  \begin{colorblock}{8. References and links}
    [1] A. Author. \emph{Example reference on research methods}.
    Example Journal, 2024.
    \par\medskip
    [2] B. Researcher. \emph{Example reference on evaluation}.
    Example Conference, 2025.
    \par\medskip
    [3] C. Collaborator. \emph{Example reference on reproducibility}.
    Example Workshop, 2026.
    \par\bigskip
    \textbf{Project:} \url{https://example.org/project}
    \par\medskip
    \textbf{Data:} \url{https://example.org/data}
  \end{colorblock}
\end{frame}
\end{document}
